\section{配分函数把能量变成全局概率}
\label{sec:13-Machine-Learning/08-Sampling-and-Inference/06}

产线上的晶圆缺陷检测用了一个打分网络：输入一张图，输出一个标量分数，分数越低越像正常品。它排序排得很好，挑出的可疑片子确实有问题。麻烦出在质量部门要的那句话上——``这张片子异常的概率是多少''。分数是相对的，你说 A 比 B 更可疑没问题，可要把 \(-3.7\) 这个分数翻译成一个概率，就得知道所有可能的图各自得几分。这不是接口设计的疏忽，是能量与概率之间隔着的那道坎。

能量模型给每个状态 \(x\) 指定一个标量 \(E_\theta(x)\)，再定义

\[
p_\theta(x)=\frac{\exp[-E_\theta(x)]}{Z(\theta)},
\qquad
Z(\theta)=\int\exp[-E_\theta(x)]\,dx
\]

（离散状态把积分换成求和）。能量只编码相对偏好，配分函数 \(Z(\theta)\) 把总质量归一化之后才有概率可言。顺带可以看出能量的绝对零点不可辨认：给 \(E_\theta\) 加一个与 \(x\) 无关的常数，分子与 \(Z\) 同乘一个因子，分布不动。

\subsection{局部能量便宜，配分函数贵在全局耦合}

\(E_\theta(x)\) 可以由一张稀疏图或一个神经网络在毫秒内算出，而 \(Z\) 要遍历整个状态空间。\(d\) 个二值变量就是 \(2^d\) 项求和，\(d=100\) 已经超出任何枚举手段。知道一张图的能量不等于知道它的概率，因为概率取决于它的分数相对于所有没见过的配置排在什么位置。

\begin{center}
\begin{tikzpicture}[>=stealth]
\begin{scope}[shift={(0,3.1)}]
\draw[->] (-3.9,0) -- (4.1,0) node[right]{\(x\)};
\draw[->] (-3.9,-0.1) -- (-3.9,1.9) node[above]{\(E_\theta(x)\)};
\draw[very thick] plot[smooth] coordinates
  {(-3.6,1.30) (-3.0,0.95) (-2.4,0.50) (-1.8,0.28) (-1.2,0.62)
   (-0.6,1.15) (0,1.45) (0.6,1.20) (1.1,0.90) (1.5,0.75)
   (2.0,0.95) (2.6,1.25) (3.2,1.50) (3.8,1.62)};
\fill (-1.8,0.28) circle (1.6pt) node[below=3pt] {\(x\)};
\fill (1.5,0.75) circle (1.6pt) node[below=3pt] {\(x'\)};
\draw[gray,dashed] (-1.8,0.28) -- (3.2,0.28);
\draw[gray,dashed] (1.5,0.75) -- (3.2,0.75);
\draw[<->,gray] (3.2,0.28) -- (3.2,0.75);
\node[gray,right] at (3.25,0.52) {\(\Delta E\)};
\end{scope}
\draw[->] (-3.9,0) -- (4.1,0) node[right]{\(x\)};
\draw[->] (-3.9,-0.1) -- (-3.9,1.9) node[above]{\(e^{-E_\theta(x)}\)};
\fill[gray!25] (-3.6,0) -- plot[smooth] coordinates
  {(-3.6,0.27) (-3.0,0.39) (-2.4,0.61) (-1.8,0.76) (-1.2,0.54)
   (-0.6,0.32) (0,0.24) (0.6,0.30) (1.1,0.41) (1.5,0.47)
   (2.0,0.39) (2.6,0.29) (3.2,0.22) (3.8,0.20)} -- (3.8,0) -- cycle;
\draw[very thick] plot[smooth] coordinates
  {(-3.6,0.27) (-3.0,0.39) (-2.4,0.61) (-1.8,0.76) (-1.2,0.54)
   (-0.6,0.32) (0,0.24) (0.6,0.30) (1.1,0.41) (1.5,0.47)
   (2.0,0.39) (2.6,0.29) (3.2,0.22) (3.8,0.20)};
\node at (0.4,1.35) {阴影总面积 \(=Z(\theta)\)：要遍历全部状态才算得出};
\end{tikzpicture}
\end{center}

\emph{两点的相对概率只看上图两个坑的深度差；任何一点的绝对概率都要除以下图那整块面积。}

图里那道分界解释了三件常被混为一谈的任务。比较两个状态的相对可能性时，\(Z\) 在比值里约掉，

\[
\frac{p_\theta(x)}{p_\theta(x')}=\exp\bigl\{-[E_\theta(x)-E_\theta(x')]\bigr\},
\]

排序、检索、结构化预测里的候选打分都落在这一类，缺陷检测的排序能用正是因为它只要这个比值。从模型采样时不必显式求出 \(Z\)——本单元第二节的接受率、第五节的权重归一化都把它消掉了——但你得造出一条真能探索 \(p_\theta\) 的链，困难从求和转成了混合速度与 Monte Carlo 误差。到了报告绝对概率、跨输入比较似然、或者按最大似然学参数的时候，\(Z\) 及其梯度会原样回来。``不知道归一化常数也能跑 MCMC''从来不等于``配分函数不重要''。

\subsection{似然梯度的负相要求你从模型里采样}

单个观测的对数似然是 \(\log p_\theta(x)=-E_\theta(x)-\log Z(\theta)\)。对第二项求导，把导数搬进积分再凑出密度，得到

\[
\nabla_\theta\log Z(\theta)=-\E_{p_\theta}\bigl[\nabla_\theta E_\theta(X)\bigr],
\]

于是

\[
\nabla_\theta\log p_\theta(x)
=-\nabla_\theta E_\theta(x)+\E_{p_\theta}\bigl[\nabla_\theta E_\theta(X)\bigr].
\]

这一步的条件对账很短：需要能在积分号下求导（被积函数与其梯度可控），需要 \(Z\) 有限。两条不满足时整个模型本身就没有定义。

结果的形状值得盯着看一会儿。第一项把数据附近的能量往下压，称正相；第二项把模型当前偏爱的区域往上抬，称负相。只做正相会让所有训练样本的能量一路下降，而没有任何东西约束别处，最后得到的是一个处处低能量、相对概率全错的模型。负相要求对模型分布取期望——也就是要从 \(p_\theta\) 采样——本单元前几节的全部工具在这里被重新调用一遍，训练的瓶颈就落在这里。

能量线性参数化时，\(E_\theta(x)=-\theta^{\trans}T(x)\)，梯度化简为

\[
\nabla_\theta\log p_\theta(x)=T(x)-\E_{p_\theta}[T(X)],
\]

即数据统计量与模型统计量之差，梯度为零就是矩匹配。这与 \hyperref[sec:13-Machine-Learning/06-Probabilistic-Models/06]{``指数族把归一化、矩与似然连接起来''} 中的结论、以及 \hyperref[sec:07-Information-Theory/01-Information-and-Entropy/04]{``最大熵原理：从约束到分布''} 里由约束导出分布的路径，是同一个结构的三种说法。差别只在模型期望能不能借图结构精确算出来。

\subsection{RBM 的条件分布简单，不代表边缘分布简单}

二值受限 Boltzmann 机把变量分成可见层 \(v\) 与隐层 \(h\)，能量为

\[
E(v,h)=-a^{\trans}v-b^{\trans}h-h^{\trans}Wv.
\]

同层无边，于是给定 \(v\) 时各隐单元条件独立，\(p(h_i=1\given v)=\sigma(b_i+W_{i:}v)\)；给定 \(h\) 时各可见单元条件独立，\(p(v_j=1\given h)=\sigma(a_j+W_{:j}^{\trans}h)\)。整层可以并行采样，交替 Gibbs 一步极便宜。便宜来自二部图结构，它不传染给边缘分布 \(p(v)\)，更不传染给 \(Z\)。正相里的 \(p(h\given v)\) 可以精确求期望，负相仍旧只能靠链。

对比散度从数据样本出发只走几步 Gibbs，用重构样本冒充模型样本。成本降下来了，短链的分布却通常不是当前模型的平稳分布，更新方向不是精确的极大似然梯度。重构图像看着像原图，只说明链在数据附近几步走不远——这恰恰是混合慢的表现，被读成了训练成功的证据。持久对比散度跨批次保留链状态，离平稳更近一些，仍然逃不出多峰能量面上混合缓慢这条老问题，而这正是第二节那张被困住的轨道图讲的事。

\subsection{归一化的范围决定了模型能回答什么问题}

能量函数是一种参数化方式，不绑定任何特定网络；决定它角色的是在什么范围上归一化。无条件能量模型在全部 \(x\) 上归一化，目标是 \(p(x)\)，\(Z\) 最难。条件能量模型对固定输入下的候选 \(y\) 归一化，得到 \(p(y\given x)\)，\hyperref[sec:13-Machine-Learning/07-Graphical-and-Sequential-Models/04]{``CRF 对整条标签序列做条件归一化''} 是结构化的典型，配分函数靠链式动态规划精确求得。softmax 分类器则是把有限个 logits 归一化，配分函数只有几十项，一次求和就算完——它也是条件能量模型，只是 \(Z\) 小到没人愿意提它。同一个公式，三种归一化范围，难度差出天壤。

\hyperref[sec:14-Deep-Learning/06-Variational-Autoencoders/01]{``VAE 把不可积后验变成可训练下界''} 走的是另一条路：显式写出联合生成过程 \(p_\theta(x,z)=p(z)p_\theta(x\given z)\)，用 \(q_\phi(z\given x)\) 近似隐变量后验，难点在对 \(z\) 的边缘化，ELBO 给出可训练的下界。能量模型的难点是全空间归一化与从模型采样。两者都用神经网络、都可能用变分近似、都能生成样本，但绕开的障碍不是同一个，不能因为输出都像数据就当成一回事。

回到质量部门那个要求。只排序候选时，未归一化的能量已经够用；要报出概率、要跨输入比较似然、要生成样本或做校准，归一化与采样质量就跳不过去。选能量模型的正当理由是相对相容性容易表达，不是把 \(Z\) 当成一个可以从公式里省略的常数。

判断近似训练可不可信，目标函数下降只是第一层。还要看独立链之间是否混合、模型样本有没有覆盖数据的多个模态、不同初值下的统计量是否一致；用短链近似时要写明链长与初始化，并检查加长链会不会改变结论；在能枚举的小状态空间上把 \(Z\) 算精确，用它核对梯度实现。最危险的失败形态是局部重构漂亮而全局分布错误：模型在每个训练点附近挖出很深的能量谷，却在谷与谷之间、在远离数据的地方堆了大量没人发现的概率质量。

整个单元的六节都绕着同一个障碍打转。目标量写得出而算不出时，摆在面前的路只有两条：采样，用样本的经验分布逼近它，要认下 Monte Carlo 误差、自相关和跨模态失败；优化，用一族简单分布逼近它，要认下族外的形状永远表达不出、反向 KL 系统性地压窄不确定性。维数高、模态少、算力紧、精度要求适中，走优化那条路；模态多、尾部要准、结论要经得起审计，且算力允许，走采样那条路并把多链诊断做全。两条路都有各自的报警器——\(\widehat R\) 与有效样本量、权重的有效样本量、族容量的敏感性检查——它们全都只能证伪，不能证实。愿意在报告里写下自己用的是哪条路、报警器读数如何、以及哪些结论会随之改变的人，才算真的把这个障碍处理过一遍。
